% Source: physical pp. 86--91 (printed pp. 63--68). Coverage: \S16.1 through the paragraph before \S16.2; Eqs. (16.1.1)--(16.1.21), all unnumbered displays, citations [1,2], and one footnote. Uncertainties: none.
\subsection{The Quantum Effective Action}
\label{sec:16.1}

Consider a quantum field theory with action $I[\phi]$, and suppose we `turn
on' a set of classical currents $J_r(x)$ coupled to the fields $\phi^r(x)$ of the
theory. The complete vacuum amplitude in the presence of these
currents is then
\begin{align}
Z[J]
&\equiv{\OutBra{\Omega}\InKet{\Omega}}_J \notag\\
&=\int\!\left[\prod_{s,y}d\phi^s(y)\right]
 \exp\!\left(iI[\phi]+i\int d^4x\,\phi^r(x)J_r(x)+\epsilon\text{ terms}\right).
\label{eq:16.1.1}
\end{align}
(The fields $\phi^r(x)$ need not be scalars. They might even be fermionic,
though until Section 16.4 we shall not bother to keep track of the signs
that would appear in that case.) The Feynman rules for calculating $Z[J]$
are just the same as for calculating the vacuum--vacuum amplitude $Z[0]$ in
the absence of the external current, except that the Feynman diagrams now
contain vertices of a new kind, to which a \emph{single} $\phi^r$-line is attached. Such
a vertex labelled with a coordinate $x$ contributes a `coupling' factor $iJ_r(x)$
to the integrand of the position-space Feynman amplitude. Equivalently,
we could say that in the expansion of $Z[J]$ in powers of $J$, the coefficient
of the term proportional to $iJ_r(x)\,iJ_s(y)\cdots$ is just the sum of diagrams
with external lines (including propagators) corresponding to the fields
$\phi^r(x)$, $\phi^s(y)$, etc. In particular, the first derivative gives the vacuum
matrix element of the quantum mechanical operator $\Phi^r(x)$ corresponding
to $\phi^r(x)$:
\begin{align}
\left[\frac{\delta}{\delta J_r(y)}Z[J]\right]_{J=0}
&=i\int\!\left[\prod_{r,x}d\phi^r(x)\right]
 \phi^r(y)\exp\{iI[\phi]+\epsilon\text{ terms}\} \notag\\
&=i\,\OutBra{\Omega}\Phi^r(y){\InKet{\Omega}}_{J=0}.
\label{eq:16.1.2}
\end{align}
We sometimes work with functionals $Z[J]$ defined by \eqref{eq:16.1.1} where the
$\phi^r(x)$ are not elementary fields (that is, fields appearing in the action) but
products of such fields. Where $\phi^r(x)$ is a product of $N$ elementary fields,
the new vertices in the Feynman rules for $Z[J]$ have $N$ lines attached.
Some of the results of this chapter (including Eq.\ \eqref{eq:16.1.2}) apply in this
case, but where Feynman diagrams are involved, it will be tacitly assumed
that the $\phi^r(x)$ are elementary.

Now, $Z[J]$ is given by the sum of \emph{all} vacuum--vacuum amplitudes in
the presence of the current $J$, including disconnected as well as connected
diagrams, but not counting as different those diagrams that differ only
by a permutation of vertices in the same or different connected subdiagrams.
A general diagram that consists of $N$ connected components will
contribute to $Z[J]$ a term equal to the product of the contributions of
these components, divided by the number $N!$ of permutations of vertices
that merely permute all the vertices in one connected component with all
the vertices in another.\footnote{The contribution of a Feynman diagram with $N$ connected components containing
$n_1,n_2,\ldots,n_N$ vertices will be proportional to a factor $1/(n_1+\cdots+n_N)!$ from the Dyson
expansion, and a factor $(n_1+\cdots+n_N)!/N!$ equal to the number of permutations of these
vertices, counting as identical those permutations that merely permute all the vertices
in one component with all the vertices in another.} Hence, the sum of all graphs is
\begin{equation}
Z[J]=\sum_{N=0}^{\infty}\frac{1}{N!}\bigl(iW[J]\bigr)^N
=\exp\bigl(iW[J]\bigr),
\label{eq:16.1.3}
\end{equation}
where $iW[J]$ is the sum of all \emph{connected} vacuum--vacuum amplitudes, again
not counting as different those diagrams that differ only by a permutation
of vertices.

For many purposes, it will be useful to go one step further, and in place
of $W[J]$ work with the sum of all connected \emph{one-particle-irreducible} graphs.
(A one-particle-irreducible graph is one that cannot be disconnected by
cutting through any one internal line.) We can give a formal expression
for this sum as follows. First, define $\phi_J^r(x)$ as the vacuum expectation
value of the operator $\Phi^r(x)$ in the presence of the current $J$:
\begin{equation}
\phi_J^r(x)\equiv
\frac{\OutBra{\Omega}\Phi^r(x){\InKet{\Omega}}_J}
{{\OutBra{\Omega}\InKet{\Omega}}_J}
=-\frac{i}{Z[J]}\frac{\delta}{\delta J_r(x)}Z[J]
\label{eq:16.1.4}
\end{equation}
or in terms of the sum of connected graphs
\begin{equation}
\phi_J^r(x)=\frac{\delta}{\delta J_r(x)}W[J].
\label{eq:16.1.5}
\end{equation}
This formula can be inverted. Define $J_{\phi r}(x)$ as the current for which
\eqref{eq:16.1.4} has a prescribed value $\phi^r(x)$:
\[
\phi_J^r(x)=\phi^r(x)\quad\text{if}\quad J_r(x)=J_{\phi r}(x).
\]
The \emph{quantum effective action}~\hyperref[ch16-ref-1]{[1]} $\Gamma[\phi]$ is defined (as a functional of $\phi$, not $J$)
by the Legendre transformation
\begin{equation}
\Gamma[\phi]\equiv-\int d^4x\,\phi^r(x)J_{\phi r}(x)+W[J_\phi].
\label{eq:16.1.6}
\end{equation}
We will soon show that $\Gamma[\phi]$ is the sum of all connected
one-particle-irreducible graphs in the presence of the current $J_\phi$. However, let us first
take a look at another aspect of its physical significance.

Note that the variational derivative of $\Gamma[\phi]$ is
\begin{align*}
\frac{\delta\Gamma[\phi]}{\delta\phi^s(y)}
={}&-\int d^4x\,\phi^r(x)\frac{\delta J_{\phi r}(x)}{\delta\phi^s(y)}-J_{\phi s}(y)\\
&+\int d^4x\left[\frac{\delta W[J]}{\delta J_r(x)}\right]_{J=J_\phi}
\frac{\delta J_{\phi r}(x)}{\delta\phi^s(y)}
\end{align*}
or, using Eq.\ \eqref{eq:16.1.5},
\begin{equation}
\frac{\delta\Gamma[\phi]}{\delta\phi^s(y)}=-J_{\phi s}(y).
\label{eq:16.1.7}
\end{equation}
Thus $\Gamma[\phi]$ is the `effective action', in the sense that the possible values for
the external fields $\phi^s(y)$ in the \emph{absence} of a current $J$ are given by the
stationary `points' of $\Gamma$:
\begin{equation}
\frac{\delta\Gamma[\phi]}{\delta\phi^s(y)}=0\quad\text{for}\quad J=0.
\label{eq:16.1.8}
\end{equation}
This may be compared with the classical field equations, which just
require that the actual action $I[\phi]$ be stationary. Hence Eq.\ \eqref{eq:16.1.8} may
be regarded as the equation of motion for the external field $\phi$, taking
quantum corrections into account.

Not only does $\Gamma[\phi]$ provide the quantum-corrected field equations; it
is an effective action in the sense that $iW[J]$ may be calculated as a sum
of connected \emph{tree} graphs for the vacuum--vacuum amplitude, with vertices
calculated as if the action were $\Gamma[\phi]$ instead of $I[\phi]$. By a tree graph is
meant one that becomes disconnected if we cut any internal line. All the
effects of loop diagrams are taken into account by using $\Gamma[\phi]$ in place of
$I[\phi]$.

To see this~\hyperref[ch16-ref-2]{[2]}, let us consider the quantity $W_\Gamma[J,g]$ that we would get
instead of $W[J]$, if we used an action $g^{-1}\Gamma[\phi]$ in place of $I[\phi]$:
\begin{align}
\exp\{iW_\Gamma[J,g]\}
\equiv{}&\int\prod_{r,x}d\phi^r(x)\,
\exp\!\left\{ig^{-1}\left[\Gamma[\phi]+\int d^4x\,\phi^r(x)J_r(x)\right]
+\epsilon\text{ terms}\right\},
\label{eq:16.1.9}
\end{align}
with arbitrary constant $g$. The propagator here is the inverse of the
coefficient of the term in $g^{-1}\Gamma[\phi]$ quadratic in $\phi$, and is hence proportional
to $g$, while all vertices make a contribution proportional to $1/g$, so a graph
with $V$ vertices (including those produced by the current $J$) and $I$ internal
lines (including those attached to the $J$ vertices) is proportional to $g^{I-V}$.
For any connected graph, the number of loops is $L=I-V+1$, so the
$L$-loop term in $W_\Gamma[J,g]$ has the $g$ dependence
\begin{equation}
\bigl(W_\Gamma[J,g]\bigr)_{L\text{ loops}}\propto g^{L-1}.
\label{eq:16.1.10}
\end{equation}
Equivalently, we may write (at least formally)
\begin{equation}
W_\Gamma[J,g]=\sum_{L=0}^{\infty}g^{L-1}W_\Gamma^{(L)}[J],
\label{eq:16.1.11}
\end{equation}
where (as can be seen by setting $g=1$), the quantity $W_\Gamma^{(L)}[J]$ is the $L$-loop
contribution to the connected vacuum amplitude $W[J,1]$ that we would
obtain if we used $\Gamma[\phi]$ (without a factor $g$) in place of the action $I[\phi]$.

Now, we are specially interested here in the sum of tree graphs, those
without loops, calculated with vertices and propagators calculated as if
the action were $\Gamma[\phi]$ instead of $I[\phi]$. In our present notation, this is
$W_\Gamma^{(0)}[J]$. In order to isolate the $L=0$ term in Eq.\ \eqref{eq:16.1.11}, consider the
limit $g\to0$. In this limit, the path integral \eqref{eq:16.1.9} is dominated by the
point of stationary phase,
\begin{equation}
\exp\{iW_\Gamma[J,g]\}\propto
\exp\!\left\{ig^{-1}\left[\Gamma[\phi_J]+\int d^4x\,\phi_J^r(x)J_r(x)\right]\right\},
\label{eq:16.1.12}
\end{equation}
where, because of its definition as the field produced by the current $J$, the
field $\phi_J$ is the stationary point of the exponent, in the sense that
\begin{equation}
\left.\frac{\delta\Gamma[\phi]}{\delta\phi^r(x)}\right|_{\phi=\phi_J}=-J_r(x).
\label{eq:16.1.13}
\end{equation}
The proportionality factor in Eq.\ \eqref{eq:16.1.12} is in general a functional of $J$,
but it is a power series in $g$ starting with terms of order $g^0$. Hence taking
the logarithm of both sides and isolating the terms of order $g^{-1}$ gives
\begin{equation}
W_\Gamma^{(0)}[J]=\Gamma[\phi_J]+\int d^4x\,\phi_J^r(x)J_r(x).
\label{eq:16.1.14}
\end{equation}
By setting $\phi=\phi_J$ in Eq.\ \eqref{eq:16.1.6}, we see that the right-hand side of
Eq.\ \eqref{eq:16.1.14} is just $W[J]$:
\begin{equation}
W_\Gamma^{(0)}[J]=W[J].
\label{eq:16.1.15}
\end{equation}
To recapitulate, this says that $W[J]$ may be calculated by using $\Gamma[\phi]$ in
place of $I[\phi]$ (the subscript $\Gamma$) and keeping only tree (0-loop) graphs:
\begin{equation}
iW[J]=\int_{\substack{\mathrm{CONNECTED}\\\mathrm{TREE}}}
\left[\prod_{r,x}d\phi^r(x)\right]
\exp\!\left\{i\Gamma[\phi]+i\sum_r\int\phi^r(x)J_r(x)\,d^4x\right\}.
\label{eq:16.1.16}
\end{equation}
Now, any connected graph for $iW[J]$ can be regarded as a tree, whose
vertices consist of one-particle-irreducible subgraphs. Thus in order for
Eq.\ \eqref{eq:16.1.16} to be correct, $i\Gamma[\phi]$ must be the sum of all
one-particle-irreducible connected graphs with arbitrary numbers of external lines,
each external line corresponding to a factor $\phi$ rather than a propagator
or wave function. For this reason, the coefficients in an expansion of $\Gamma[\phi]$
in powers of fields and their derivatives around some fixed field $\phi_0$ may
be regarded as \emph{renormalized} coupling constants, with the renormalization
`point' specified by $\phi_0$ rather than by some set of momenta.

Equivalently, $i\Gamma[\phi_0]$ for some fixed field $\phi_0^r(x)$ may be expressed as the
sum of one-particle-irreducible graphs for the vacuum--vacuum amplitude,
calculated with a shifted action $I[\phi+\phi_0]$:
\begin{equation}
i\Gamma[\phi_0]=\int_{\substack{\mathrm{1PI}\\\mathrm{CONNECTED}}}
\left[\prod_{r,x}d\phi^r(x)\right]\exp\{iI[\phi+\phi_0]\}.
\label{eq:16.1.17}
\end{equation}
This is because any place where $\phi_0$ appears in any of the vertices or
propagators within the one-particle-irreducible graphs in Eq.\ \eqref{eq:16.1.17} is
also a place where an external $\phi$-line could be attached. (The restriction
to one-particle-irreducible graphs plays an essential role in Eq.\ \eqref{eq:16.1.17};
without this restriction we could shift the variable of integration, yielding
an integral that would be manifestly independent of $\phi_0$.) In place of
Eq.\ \eqref{eq:16.1.17}, it is often convenient to write
\begin{equation}
\exp\{i\Gamma[\phi_0]\}=\int_{\mathrm{1PI}}
\left[\prod_{r,x}d\phi^r(x)\right]\exp\{iI[\phi+\phi_0]\},
\label{eq:16.1.18}
\end{equation}
in which we evaluate the path integral including all graphs, connected or
not, in which each connected component is one-particle-irreducible.

\begin{center}
$*\quad*\quad*$
\end{center}

This formalism provides a simple method of summing tree graphs. As
one example, consider the relation between the complete two-point function
$\Delta^{rx,sy}$ and its one-particle-irreducible part $\Pi_{rx,sy}$. From Eqs.\ \eqref{eq:16.1.5}
and \eqref{eq:16.1.7}, we find
\begin{equation}
\Delta^{rx,sy}\equiv\frac{\delta^2W[J]}{\delta J_r(x)\,\delta J_s(y)}
=\frac{\delta\phi_J^r(x)}{\delta J_s(y)},
\label{eq:16.1.19}
\end{equation}
\begin{equation}
\Pi_{rx,sy}\equiv\frac{\delta^2\Gamma[\phi]}{\delta\phi^r(x)\,\delta\phi^s(y)}
=-\frac{\delta J_{\phi r}(x)}{\delta\phi^s(y)}.
\label{eq:16.1.20}
\end{equation}
It follows immediately that the `matrices' $\Delta$ and $\Pi$ are related by
\begin{equation}
\Delta=-\Pi^{-1}.
\label{eq:16.1.21}
\end{equation}
This is the counterpart of the familiar relation (10.3.15) between propagators
and self-energy parts, with the extra term $q^2+m^2$ in the denominator
in Eq.\ (10.3.15) representing the zeroth-order term in the
one-particle-irreducible two-point function.
